\documentclass[11pt,a4paper]{report}
\usepackage[utf8]{inputenc}
\usepackage[T1]{fontenc}
\usepackage[italian]{babel}
\usepackage{amsmath, amssymb, bm}
\usepackage{graphicx}
\usepackage{hyperref}
\usepackage{listings}
\usepackage{xcolor}
\usepackage{geometry}
\usepackage{booktabs}
\usepackage{caption}
\usepackage{mdframed}
\geometry{a4paper, margin=2.5cm}

\definecolor{codegreen}{rgb}{0,0.55,0}
\definecolor{codegray}{rgb}{0.5,0.5,0.5}
\definecolor{codepurple}{rgb}{0.55,0,0.75}
\definecolor{backcolour}{rgb}{0.97,0.97,0.94}
\definecolor{theoryblue}{rgb}{0.90,0.95,1.0}

\lstdefinestyle{mystyle}{
    backgroundcolor=\color{backcolour},
    commentstyle=\color{codegreen}\itshape,
    keywordstyle=\color{magenta}\bfseries,
    numberstyle=\tiny\color{codegray},
    stringstyle=\color{codepurple},
    basicstyle=\ttfamily\small,
    breakatwhitespace=false,
    breaklines=true,
    captionpos=b,
    keepspaces=true,
    numbers=left,
    numbersep=5pt,
    showspaces=false,
    showstringspaces=false,
    showtabs=false,
    tabsize=4,
    frame=single,
    rulecolor=\color{gray!40}
}
\lstset{style=mystyle}

\newmdenv[backgroundcolor=theoryblue, linecolor=blue!40, roundcorner=4pt]{theoryframe}

\title{
    \vspace{2cm}
    \textbf{\Huge Wheeled Biped Robot (RoTino)}\\
    \vspace{0.5cm}
    \Large Manuale Tecnico Completo:\\
    Architettura del Codice e Fondamenti Teorici\\
    \vspace{0.3cm}
    \normalsize Analisi file per file con collegamento alla teoria del controllo
    \vspace{2cm}
}
\author{Analisi Tecnica Approfondita}
\date{\today}

\begin{document}

\maketitle
\tableofcontents
\newpage

%=============================================================
\chapter{Introduzione e Architettura del Progetto}
%=============================================================

\section{Obiettivo del Documento}
Questo manuale ha lo scopo di collegare in modo esplicito e rigoroso ogni
singolo file sorgente del progetto RoTino alla teoria del controllo e della
robotica da cui trae origine. Per ciascun file viene presentata prima una
descrizione del ruolo architetturale, poi le righe di codice piu' significative
con la corrispondente derivazione matematica.

\section{Albero Completo dei File}

Di seguito e' riportata la struttura completa del workspace ROS~2 Humble:

\begin{verbatim}
rotino_ws/
|-- src/
|   |-- rotino_description/
|   |   |-- rotino_description/
|   |   |   |-- model.py              <-- Parametri VL-WIP dall'URDF
|   |   |   |-- kinematics.py         <-- FK analitica e CoM globale
|   |   |   `-- planar_trajectory.py  <-- Legge oraria S-curve quintica
|   |   `-- urdf/
|   |       `-- rotino.urdf.xacro     <-- Geometria, masse, giunti
|   |
|   |-- rotino_pid/
|   |   `-- rotino_pid/
|   |       `-- controller.py         <-- Bilanciamento PD + jump FSM
|   |
|   |-- rotino_smc/
|   |   `-- rotino_smc/
|   |       |-- controller.py         <-- Super-Twisting + Task-Space VMC
|   |       `-- estimator.py          <-- Filtro di Kalman lineare
|   |
|   |-- rotino_mpc/
|   |   `-- rotino_mpc/
|   |       |-- controller.py         <-- Nodo ROS 2, loop di controllo
|   |       `-- solvers.py            <-- QP FISTA, LQR, condensazione
|   |
|   `-- rotino_benchmark/
|       `-- rotino_benchmark/
|           |-- campaign.py           <-- Orchestrazione run sequenziali
|           |-- logger.py             <-- Logging sincrono ad alta freq.
|           |-- plot.py               <-- 9 pannelli di analisi
|           `-- compare.py            <-- Confronto tra run
`-- docs/
    `-- Architettura_e_Codice.tex     <- Questo documento
\end{verbatim}

%=============================================================
\chapter{rotino\_description: Il Modello Dinamico del Robot}
%=============================================================

\section{Ruolo del Pacchetto}
Questo pacchetto ha due compiti fondamentali: (1) definire la geometria
e le inerzie del robot tramite URDF/Xacro, e (2) fornire a tutti gli altri
pacchetti una libreria Python che estrae automaticamente i parametri
dinamici dal modello URDF, evitando valori scritti a mano che potrebbero
discostarsi dalla simulazione.

%-------------------------------------------------------------
\section{File: kinematics.py}
%-------------------------------------------------------------

\subsection{Ruolo}
Implementa la \textbf{cinematica diretta (Forward Kinematics)} del robot
traversando l'albero URDF a partire dalla radice, senza dipendenze da
librerie esterne come Pinocchio o MoveIt. Calcola il centro di massa
globale tramite media ponderata delle masse.

\subsection{Riga 9--17: Conversione RPY a matrice di rotazione}

\begin{theoryframe}
\textbf{Teoria:} Una rototraslazione rigida nello spazio 3D si rappresenta con
la matrice di rotazione $\mathbf{R} \in SO(3)$. Per la convenzione RPY (Roll,
Pitch, Yaw) usata da URDF e ROS, la rotazione composta e' il prodotto:
$\mathbf{R} = R_z(\psi) R_y(\phi) R_x(\theta)$, la cui forma esplicita e' la
classica matrice di rotazione di Eulero ZYX.
\end{theoryframe}

\begin{lstlisting}[language=Python, caption=kinematics.py rr.\ 9-17: Matrice di rotazione da angoli RPY]
def rpy_to_matrix(r, p, y):
    cr, sr = np.cos(r), np.sin(r)
    cp, sp = np.cos(p), np.sin(p)
    cy, sy = np.cos(y), np.sin(y)
    # Prodotto R_z * R_y * R_x in forma esplicita (convenzione ZYX)
    return np.array([
        [cy*cp, cy*sp*sr - sy*cr, cy*sp*cr + sy*sr],
        [sy*cp, sy*sp*sr + cy*cr, sy*sp*cr - cy*sr],
        [-sp,   cp*sr,            cp*cr],
    ])
\end{lstlisting}

Questa funzione viene richiamata ogni volta che si propaga la trasformazione
di un giunto fisso dell'URDF, che specifica la sua posa relativa tramite
attributo \texttt{rpy} nella tag \texttt{<origin>}.

\subsection{Riga 20--24: Formula di Rodrigues (Asse-Angolo)}

\begin{theoryframe}
\textbf{Teoria (Formula di Rodrigues):} Data una rotazione di angolo $\alpha$
attorno all'asse unitario $\hat{\mathbf{k}}$, la matrice di rotazione e':
\[
\mathbf{R} = \mathbf{I} + \sin\alpha\,[\hat{\mathbf{k}}]_\times + (1-\cos\alpha)\,[\hat{\mathbf{k}}]_\times^2
\]
dove $[\cdot]_\times$ e' la matrice antisimmetrica (skew-symmetric) tale che
$[\mathbf{k}]_\times \mathbf{v} = \mathbf{k} \times \mathbf{v}$.
\end{theoryframe}

\begin{lstlisting}[language=Python, caption=kinematics.py rr.\ 20-24: Formula di Rodrigues]
def axis_angle_to_matrix(axis, angle):
    k = np.asarray(axis, dtype=float)
    k = k / np.linalg.norm(k)
    # kx e' la matrice skew-symmetric [k]_x
    kx = np.array([[0.0, -k[2], k[1]],
                   [k[2],  0.0, -k[0]],
                   [-k[1], k[0],  0.0]])
    # Formula di Rodrigues esatta
    return np.eye(3) + np.sin(angle) * kx + (1.0 - np.cos(angle)) * (kx @ kx)
\end{lstlisting}

Questa e' la formula canonica per costruire la matrice di rotazione $R \in SO(3)$
attorno a un asse arbitrario. Viene usata in \texttt{link\_frames} per ogni
giunto revoluto, dove l'asse e' quello specificato dall'URDF e l'angolo e'
la variabile di giunto $q_i$.

\subsection{Riga 54--71: Cinematica Diretta ad Albero}

\begin{theoryframe}
\textbf{Teoria (FK per strutture ad albero):} La cinematica diretta propaga
la trasformazione omogenea ${}^0T_n = {}^0T_1 \cdot {}^1T_2 \cdots {}^{n-1}T_n$
dal nodo radice verso le foglie. In forma compatta, per ogni link figlio:
$p_{child} = p_{parent} + R_{parent} \cdot p_{joint}$,
$R_{child} = R_{parent} \cdot R_{joint} \cdot R_q(q_i)$.
\end{theoryframe}

\begin{lstlisting}[language=Python, caption=kinematics.py rr.\ 54-71: FK iterativa via stack DFS]
def link_frames(self, base_pos, base_rot, joint_positions):
    frames = {}
    stack = [(self.root, base_rot, base_pos)]
    while stack:
        name, rot, pos = stack.pop()
        frames[name] = (rot, pos)
        for joint in self.child_joints.get(name, []):
            origin = joint.origin
            xyz = np.array(origin.xyz if origin.xyz else [0., 0., 0.])
            rpy = origin.rpy if origin.rpy else [0., 0., 0.]
            # p_child = p_parent + R_parent * p_joint_origin
            child_pos = pos + rot @ xyz
            # R_child = R_parent * R_joint_origin
            child_rot = rot @ rpy_to_matrix(*rpy)
            if joint.type in ('revolute', 'continuous'):
                # Applica la rotazione del giunto con la formula di Rodrigues
                child_rot = child_rot @ axis_angle_to_matrix(
                    joint.axis, joint_positions.get(joint.name, 0.0))
            stack.append((joint.child, child_rot, child_pos))
    return frames
\end{lstlisting}

L'approccio DFS (Depth First Search) tramite stack esplicito e' equivalente
alla ricorsione ma evita l'overhead dello stack Python per alberi profondi.

%-------------------------------------------------------------
\section{File: model.py}
%-------------------------------------------------------------

\subsection{Ruolo}
Questo e' il nucleo analitico del progetto. Implementa il modello
\textbf{VL-WIP (Variable-Length Wheeled Inverted Pendulum)} di Cui et al.,
\emph{Micromachines 2022, 13, 747}, calcolando automaticamente tutti i
parametri dinamici dalla geometria URDF. Ogni calcolo ha un riferimento
esplicito all'equazione del paper.

\subsection{Riga 28--48: Struttura dati WBRParams (Tabella 1 del Paper)}

\begin{lstlisting}[language=Python, caption=model.py rr.\ 28-48: Struttura parametri VL-WIP]
@dataclass
class WBRParams:
    m_w: float   # massa di una ruota [kg]             -> Tab.1: m_w
    I_w: float   # inerzia di spin della ruota [kg m2] -> Tab.1: I_w
    r:   float   # raggio della ruota [m]              -> Tab.1: r
    d:   float   # carreggiata [m]                     -> Tab.1: d
    m_b: float   # massa corpo superiore [kg]          -> Tab.1: m_b
    m_1: float   # massa dello stinco [kg]             -> Tab.1: m_1
    m_2: float   # massa della coscia [kg]             -> Tab.1: m_2
    l_1: float   # lunghezza stinco: ginocchio->ruota  -> Tab.1: l_1
    l_2: float   # lunghezza coscia: anca->ginocchio   -> Tab.1: l_2
    L_max: float # lunghezza massima gamba = l1 + l2   -> Tab.1: L
    mu:   float  # coeff. di attrito ruota/suolo       -> Tab.1: mu
    leg_damping: float  # smorzamento viscoso giunto [Nm s/rad]
\end{lstlisting}

Questi parametri corrispondono esattamente alla Tabella 1 del paper di
riferimento. L'estrazione automatica garantisce coerenza tra simulazione e
modello analitico.

\subsection{Riga 151--178: Centroide Equivalente (Eqs. 2-3 del Paper)}

\begin{theoryframe}
\textbf{Teoria (Eqs. 2-3):} Il WBR viene modellato come un pendolo inverso
a lunghezza variabile. Il corpo superiore (gambe + torso) e' sostituito da un
punto materiale equivalente di massa $m_b$:
\[
S_C = \frac{\sum_i m_i x_i}{m_b}, \quad
Z_C = \frac{\sum_i m_i z_i}{m_b}, \quad
l = \sqrt{S_C^2 + Z_C^2}, \quad
\theta = \arctan\!\left(\frac{S_C}{Z_C}\right)
\]
dove la somma si estende a tutti i link tranne le ruote, e le coordinate sono
nel frame dell'asse ruote.
\end{theoryframe}

\begin{lstlisting}[language=Python, caption=model.py rr.\ 151-178: Centroide equivalente (Eqs. 2-3)]
def equivalent_centroid(self, hip, knee, pitch=0.0):
    R = rpy_to_matrix(0.0, pitch, 0.0)
    frames = self.kin.link_frames(np.zeros(3), R, self.leg_joints(hip, knee))
    # Punto medio dell'asse ruote
    axle = 0.5*(frames[LEFT_WHEEL_LINK][1] + frames[RIGHT_WHEEL_LINK][1])
    # CoM ponderato del corpo superiore (tutto tranne le ruote)
    masses = np.array([_link(self.robot, n).inertial.mass for n in names])
    coms = np.array([self._link_com(frames, n) for n in names])
    com = masses @ coms / masses.sum()
    # Eq. (2)-(3): distanza S_C, Z_C rispetto all'asse ruote
    S_C = float(com[0] - axle[0])
    Z_C = float(com[2] - axle[2])
    return {
        'l':     math.hypot(S_C, Z_C),      # lunghezza del pendolo equivalente
        'theta': math.atan2(S_C, Z_C),      # angolo rispetto alla verticale
        'I_y':   float(I[1, 1]),             # inerzia di pitch del corpo sup.
        'z_b':   float(-axle[2]),            # altezza asse ruote da terra
    }
\end{lstlisting}

\subsection{Riga 180--195: Coefficienti VL-WIP (Eq. 14 del Paper)}

\begin{theoryframe}
\textbf{Teoria (Eq. 14):} Linearizzando le equazioni di Lagrange del WBR
attorno all'equilibrio (robot verticale, ruote che rotolano senza slittamento),
si ottiene il modello lineare:
\[
\ddot{s} = a_1 \theta + b_1 (\tau_L + \tau_R), \qquad
\ddot{\theta} = a_2 \theta + b_2 (\tau_L + \tau_R)
\]
dove $a_2 > 0$ implica l'instabilita' del pendolo (autovalore positivo).
I coefficienti dipendono dalla lunghezza $l$ del pendolo, variabile con la postura.
\end{theoryframe}

\begin{lstlisting}[language=Python, caption=model.py rr.\ 180-195: Coefficienti VL-WIP (Eq. 14)]
def vlwip_coefficients(self, l, I_y=None, I_z=None):
    p = self.p
    if I_y is None:
        I_y = p.m_b * l * l / 3.0  # approx. del paper: corpo a bastone
    mb, mw, Iw, r, d = p.m_b, p.m_w, p.I_w, p.r, p.d
    # Denominatore comune (determinante della matrice di inerzia)
    den = (2*Iw*(I_y + mb*l*l)
           + (2*l*l*mb*mw + I_y*(mb + 2*mw))*r*r)
    return {
        'a1': -G*l*l*mb*mb*r*r / den,        # accoppiamento s <- theta
        'a2':  G*l*mb*(2*Iw+(mb+2*mw)*r*r)/den,  # instabilita' pendolo
        'b1':  r*(I_y + l*mb*(l+r)) / den,    # guadagno coppia -> s
        'b2': -(2*Iw + r*(l*mb+(mb+2*mw)*r))/den, # guadagno coppia -> theta
        'b3':  d*r / (2*I_z*r*r + d*d*(Iw+mw*r*r)),  # guadagno yaw
    }
\end{lstlisting}

Il fatto che $b_2 < 0$ e' fondamentale: nella funzione
\texttt{\_pitch\_super\_twisting}, la coppia fisica alle ruote si calcola
dividendo per $b_2$, il che inverte automaticamente il segno del segnale
di controllo.

\subsection{Riga 211--223: FK Analitica 2D delle Gambe e Jacobiano}

\begin{theoryframe}
\textbf{Teoria (Catena cinematica planare):} Per un robot a 2 link planari
(anca-ginocchio-ruota), la posizione del centro ruota nel frame base e':
$\mathbf{p}_{wheel} = \mathbf{p}_{hip} + R(\phi_{hip})\mathbf{p}_{knee|hip}
+ R(\phi_{hip}+\phi_{knee})\mathbf{p}_{wheel|knee}$.
Il Jacobiano $\mathbf{J} = \partial \mathbf{p}_{wheel}/\partial(\phi_{hip},\phi_{knee})$
si calcola derivando analiticamente rispetto agli angoli di giunto.
\end{theoryframe}

\begin{lstlisting}[language=Python, caption=model.py rr.\ 211-223: FK 2D e Jacobiano analitico]
def leg_fk(self, hip, knee):
    kx, kz = self.knee_xz    # offset ginocchio nel frame dell'anca [m]
    wx, wz = self.wheel_xz   # offset ruota nel frame del ginocchio [m]

    def rot(q, x, z):
        c, s = math.cos(q), math.sin(q)
        # Punto ruotato e sua derivata rispetto all'angolo q
        return (np.array([x*c + z*s, -x*s + z*c]),   # posizione
                np.array([-x*s + z*c, -x*c - z*s]))  # dp/dq (colonna di J)

    pk, dpk = rot(hip, kx, kz)            # posizione ginocchio e sua derivata
    pw, dpw = rot(hip + knee, wx, wz)     # posizione ruota e sua derivata
    # J[:,0] = d(p_ruota)/d(hip): ruotare l'anca muove sia ginocchio che ruota
    # J[:,1] = d(p_ruota)/d(knee): ruotare il ginocchio muove solo la ruota
    J = np.column_stack((dpk + dpw, dpw))
    return self.hip_xz + pk + pw, J
\end{lstlisting}

%-------------------------------------------------------------
\section{File: planar\_trajectory.py}
%-------------------------------------------------------------

\subsection{Ruolo}
Genera la traiettoria planare (S-curve laterale) che il robot segue durante
la navigazione in avanzamento e sterzata, con una legge oraria quintica che
garantisce accelerazione zero agli estremi (jerk minimizzato).

\subsection{Riga 50--63: Legge Oraria Quintica e Curvatura}

\begin{theoryframe}
\textbf{Teoria (Polinomio quintico per profili smooth):} Un profilo di
posizione $s(t)$ con $s(0)=\dot s(0)=\ddot s(0)=0$ e $s(T)=L$,
$\dot s(T)=\ddot s(T)=0$ (``profilo a S'') si ottiene con il polinomio quintico
normalizzato $\tau=t/T$:
\[s(\tau) = L(10\tau^3 - 15\tau^4 + 6\tau^5)\]
La velocita' $\dot s = L(30\tau^2 - 60\tau^3 + 30\tau^4)/T$ ha profilo
trapezoidale arrotondato. La curvatura della traiettoria planare e'
$\kappa = y''/(1+y'^2)^{3/2}$ e il yaw\_rate desiderato e' $\dot\psi = \kappa \cdot v$.
\end{theoryframe}

\begin{lstlisting}[language=Python, caption=planar\_trajectory.py rr.\ 50-63: Legge quintica e campionamento]
def _time_law(self, t):
    tau = min(max(t / self.duration, 0.0), 1.0)   # normalizzazione [0,1]
    s = self.length * (10*tau**3 - 15*tau**4 + 6*tau**5)     # posizione
    v = self.length * (30*tau**2 - 60*tau**3 + 30*tau**4) / self.duration # velocita'
    return s, v

def sample(self, t):
    s, v = self._time_law(t)
    # Inversione della curva lunghezza-arco (tabella precalcolata)
    x = float(np.interp(s, self.s_table, self.x_table))
    dy = self._dy(x)                           # derivata della S-curve spaziale
    heading = math.atan(dy)                    # angolo tangente alla curva
    # Curvatura della curva planare: kappa = y'' / (1 + y'^2)^(3/2)
    curvature = self._ddy(x) / (1.0 + dy*dy)**1.5
    # yaw_rate desiderato = curvatura * velocita' lineare
    return x, float(self._y(x)), heading, v, curvature * v
\end{lstlisting}

%=============================================================
\chapter{rotino\_smc: Sliding Mode Control}
%=============================================================

\section{Architettura del Controllore}
Il pacchetto implementa tre leggi di controllo simultanee:
(1) pitch Super-Twisting sulle ruote,
(2) leg Task-Space SMC (VMC esteso),
(3) yaw PD+I con feedforward di attrito.

%-------------------------------------------------------------
\section{File: estimator.py}
%-------------------------------------------------------------

\subsection{Ruolo}
Implementa il \textbf{Filtro di Kalman lineare a due stati} $[p, v]$ per
stimare posizione e velocita' del corpo superiore, fondendo le misure
dell'accelerometro IMU (aggiornamento veloce) con l'odometria delle ruote
combinata alla cinematica delle gambe (correzione lenta e accurata).

\subsection{Riga 26--37: Predizione e Correzione del Filtro di Kalman}

\begin{theoryframe}
\textbf{Teoria (Kalman Discreto):} Per il sistema $\mathbf{x}_{k+1} = F\mathbf{x}_k + Bu_k + w_k$
con $w_k \sim \mathcal{N}(0,Q)$ e osservazione $y_k = \mathbf{x}_k + v_k$
con $v_k \sim \mathcal{N}(0,R)$, le equazioni ottimali sono:\\
Predizione: $\hat x^- = F\hat x_{k-1} + Bu_{k-1}$, $P^- = FP_{k-1}F^T + Q$\\
Correzione: $K = P^-(P^- + R)^{-1}$, $\hat x = \hat x^- + K(y - \hat x^-)$,
$P = (I-K)P^-$
\end{theoryframe}

\begin{lstlisting}[language=Python, caption=estimator.py rr.\ 26-37: Filtro di Kalman (Eqs.\ 20-21 di Cui 2022)]
def predict(self, acc, dt):
    # Matrice di transizione cinematica: x=[p,v], input=accelerazione IMU
    F = np.array([[1.0, dt], [0.0, 1.0]])
    Bu = np.array([0.5*dt*dt, dt])   # vettore di ingresso (cinematica newtoniana)
    # Predizione stato: x^- = F x + B u
    self.x = F @ self.x + Bu * acc
    # Propagazione covarianza: P^- = F P F^T + Q
    # Q = q_acc * Bu Bu^T modella il rumore dell'accelerometro
    self.P = F @ self.P @ F.T + self.q_acc * np.outer(Bu, Bu) + 1e-9 * np.eye(2)

def correct(self, p_obs, v_obs, noise_scale=1.0):
    R = self.R_nom * noise_scale       # covarianza del rumore di misura
    S = self.P + R                     # covarianza dell'innovazione
    K = self.P @ np.linalg.inv(S)     # guadagno ottimale di Kalman
    self.x = self.x + K @ (np.array([p_obs, v_obs]) - self.x)  # correzione
    self.P = (np.eye(2) - K) @ self.P  # aggiornamento covarianza
\end{lstlisting}

Il termine $q_{acc} \cdot \mathbf{B}_u\mathbf{B}_u^T$ modella un rumore
sull'accelerazione (random walk in accelerazione), il che e' fisicamente
appropriato per un MEMS IMU che ha un bias lentamente variabile.

%-------------------------------------------------------------
\section{File: controller.py (SMC)}
%-------------------------------------------------------------

\subsection{Riga 576--587: Pitch Super-Twisting (Algoritmo di Levant)}

\begin{theoryframe}
\textbf{Teoria (Super-Twisting):} L'algoritmo Super-Twisting di Levant e'
uno SMC del 2$^\circ$ ordine che garantisce convergenza in tempo finito su
$\sigma=0$ con segnale di controllo \textbf{continuo} (elimina il chattering).
Data la superficie $\sigma$, la legge e':
\[u_1 = -\alpha|\sigma|^{1/2}\mathrm{sgn}(\sigma) + u_2, \qquad
\dot u_2 = -\beta\,\mathrm{sgn}(\sigma)\]
La funzione di Lyapunov $V = \alpha^2|\sigma| + \frac{1}{2}u_2^2$ dimostra
la convergenza finita.
\end{theoryframe}

\begin{lstlisting}[language=Python, caption=controller.py SMC rr.\ 576-587: Pitch Super-Twisting]
def _pitch_super_twisting(self, theta, theta_ref, l):
    c = self._vlwip(l)   # coefficienti VL-WIP alla lunghezza l corrente

    # 1. Superficie di scorrimento: su sigma=0, theta -> theta_ref esponenzialmente
    s1 = self.theta_dot + SMC_C1 * (theta - theta_ref)

    # 2. Approssimazione continua del segno (riduce chattering residuo)
    sg1 = s1 / (abs(s1) + SMC_EPS1)

    # 3. Legge astratta nu (Super-Twisting di Levant):
    #    -C1*theta_dot: termine derivativo (mantiene la superficie)
    #    -KA*sqrt(|s1|)*sgn(s1): termine radice (convergenza finita)
    #    self.W: integratore (aggiornato nel loop: W -= KB*sg1*dt)
    nu = -SMC_C1*self.theta_dot - SMC_KA*math.sqrt(abs(s1))*sg1 + self.W

    # 4. Feedback Linearization: theta_ddot = a2*theta + b2*(tau_L+tau_R)
    #    Invertendo: tau_comune = (nu - a2*theta) / b2 / 2
    #    (diviso 2: le due ruote erogano la stessa coppia)
    return 0.5 * (nu - c['a2'] * theta) / c['b2'], s1, sg1
\end{lstlisting}

\subsection{Righe di Aggiornamento dell'Integratore W nel Loop Principale}

\begin{lstlisting}[language=Python, caption=controller.py SMC: Aggiornamento dell'integratore di Levant]
# Nel loop di controllo principale, dopo _pitch_super_twisting:
# Integrazione di W_dot = -KB * sgn(s1) (integratore continuo di Levant)
# La clip su SMC_W_MAX evita il wind-up dell'integratore
self.W = float(np.clip(
    self.W - SMC_KB * sg1 * min(dt, SMC_DT_MAX),
    -SMC_W_MAX, SMC_W_MAX))
\end{lstlisting}

\subsection{Riga 555--574: Task-Space VMC (Virtual Model Control)}

\begin{theoryframe}
\textbf{Teoria (VMC + SMC nel Task Space):} Il VMC di Pratt et al. (1997)
comanda forze virtuali cartesiane e le mappa nei giunti: $\tau = \mathbf{J}^T \mathbf{F}$.
In questa implementazione avanzata, $\mathbf{F}$ include una superficie di
scorrimento cartesiana $\mathbf{s} = (\mathbf{v}_f - \mathbf{v}_d) + c(\mathbf{p}_f - \mathbf{p}_d)$:
\[\mathbf{F} = -K_{lin}\mathbf{s} - K_{sw}\,\mathrm{sat}\!\left(\frac{\mathbf{s}}{\phi}\right) + \mathbf{F}_{ff}\]
Il termine $\mathrm{sat}(\mathbf{s}/\phi)$ e' la funzione segno saturata che
approssima il controllo discontinuo SMC con un segnale continuo.
\end{theoryframe}

\begin{lstlisting}[language=Python, caption=controller.py SMC rr.\ 555-574: VMC + Sliding Mode nel Task Space]
def _leg_smc(self, legs, p_d, c, K_lin, F_ff=np.zeros(2), v_d=np.zeros(2)):
    taus = []
    for i, ((p_f, J, v_f), qd) in enumerate(zip(legs, self.leg_qd)):
        # Superficie di scorrimento nel Task Space (errore PD cartesiano)
        # p_f: posizione FK del centro ruota, p_d: setpoint cartesiano
        s = (v_f - v_d) + c * (p_f - p_d)

        # Forza virtuale composta:
        # -K_lin @ s:  termine proporzionale+derivativo (banda passante)
        # -LEG_KSW*clip(s/PHI,-1,1):  switching saturato (robustezza SMC)
        # F_ff:        feedforward peso del robot (compensa la gravita')
        F = -K_lin @ s - LEG_KSW * np.clip(s / LEG_PHI, -1.0, 1.0) + F_ff

        # Feedforward per cancellare lo smorzamento viscoso dell'URDF
        tau_ff = self.model.p.leg_damping * qd

        # LEGGE VMC: tau = J^T @ F (mappa forze cartesiane in coppie giunturali)
        taus.append(np.clip(J.T @ F + tau_ff, -LEG_TORQUE_MAX, LEG_TORQUE_MAX))

    # Riordino per il topic ROS: [hip_L, hip_R, knee_L, knee_R]
    tau = [taus[0][0], taus[1][0], taus[0][1], taus[1][1]]
    self.leg_pub.publish(Float64MultiArray(data=[float(x) for x in tau]))
    return tau
\end{lstlisting}

\subsection{Loop Esterno di Velocita' (Non-Minimum Phase)}

\begin{theoryframe}
\textbf{Teoria (Sistema a fase non minima):} Il WBR e' un sistema a fase non
minima: una coppia positiva alle ruote genera inizialmente uno spostamento del
corpo all'indietro (reazione di Newton) prima di accelerare in avanti.
Per questo, il loop PI sulla velocita' deve avere guadagni \textbf{negativi}
$K_v^P, K_v^I < 0$: un errore di velocita' positivo (troppo lento) richiede
un angolo di pitch negativo (inclinare il robot in avanti) per generare
l'accelerazione desiderata.
\end{theoryframe}

\begin{lstlisting}[language=Python, caption=controller.py SMC: Loop esterno PI (non-minimum phase)]
def _outer_velocity_loop(self, v, v_ref, dt, active):
    err = v - v_ref
    if active:
        # Integrale dell'errore di velocita' (con anti-windup)
        self.Iv = float(np.clip(self.Iv + V_KI * err * dt, -IV_MAX, IV_MAX))
    # Riferimento di pitch: negativo perche' il sistema e' a fase non minima
    theta_star = float(np.clip(V_KP * err + self.Iv, -THETA_REF_MAX, THETA_REF_MAX))
    return theta_star
\end{lstlisting}

%=============================================================
\chapter{rotino\_mpc: Model Predictive Control}
%=============================================================

\section{Architettura MPC}
Il controllore MPC separa il problema in due QP indipendenti seguendo
il modello disaccoppiato di Eq. (9) di Cui et al.: uno orizzontale
(posizione e velocita' nel piano sagittale) e uno verticale (altezza del
bacino e forza verticale alle ruote). Questo disaccoppiamento e' possibile
grazie all'approssimazione del LIP (Linear Inverted Pendulum).

%-------------------------------------------------------------
\section{File: solvers.py}
%-------------------------------------------------------------

\subsection{Riga 15--30: CARE e Guadagno LQR Ottimale}

\begin{theoryframe}
\textbf{Teoria (LQR):} Il controllore LQR minimizza il costo quadratico
$J = \int_0^\infty (x^T Q x + u^T R u) dt$ per il sistema lineare
$\dot x = Ax + Bu$. La soluzione ottimale e' $u^* = -Kx$ con
$K = R^{-1} B^T P$ dove $P$ soddisfa la CARE:
$A^T P + PA - PBR^{-1}B^T P + Q = 0$.
La soluzione si trova dall'invariante stabile dell'hamiltoniano
$H = \begin{pmatrix}A & -BR^{-1}B^T\\-Q & -A^T\end{pmatrix}$.
\end{theoryframe}

\begin{lstlisting}[language=Python, caption=solvers.py rr.\ 15-30: CARE e guadagno LQR]
def care(A, B, Q, R):
    n = A.shape[0]
    Rinv = np.linalg.inv(R)
    # Costruzione della matrice Hamiltoniana (2n x 2n)
    H = np.block([[A, -B @ Rinv @ B.T], [-Q, -A.T]])
    w, V = np.linalg.eig(H)
    # Autovettori con parte reale negativa (sottospazio stabile)
    stable = V[:, w.real < 0.0]
    # Soluzione di Riccati: P = Im_lower * Im_upper^{-1}
    P = np.real(stable[n:, :] @ np.linalg.inv(stable[:n, :]))
    return 0.5 * (P + P.T)   # simmetrizzazione numerica

def lqr_gain(A, B, Q, R):
    P = care(A, B, Q, R)
    return np.linalg.solve(R, B.T @ P)   # K = R^{-1} B^T P
\end{lstlisting}

\subsection{Riga 33--45: LQRSchedule (TV-LQR su griglia di lunghezze)}

\begin{theoryframe}
\textbf{Teoria (Time-Varying LQR):} Il guadagno LQR ottimale dipende dalla
lunghezza $l$ del pendolo, che varia con la postura delle gambe. Anziche'
risolvere la CARE a ogni passo (computazionalmente costoso), si precalcola
il guadagno su una griglia di valori di $l$ e si interpola linearmente
a runtime, realizzando un TV-LQR efficiente.
\end{theoryframe}

\begin{lstlisting}[language=Python, caption=solvers.py rr.\ 33-45: LQRSchedule con interpolazione su griglia]
class LQRSchedule:
    def __init__(self, model, samples, Q, R):
        self.l_grid = np.array([s[0] for s in samples])
        # Precalcolo offline del guadagno ottimale per ogni l nella griglia
        self.K_grid = np.array([lqr_gain(*model.vlwip_matrices(l, I_y=iy), Q, R)
                                 for l, iy in samples])

    def gain(self, l):
        l = float(np.clip(l, self.l_grid[0], self.l_grid[-1]))
        i = int(np.clip(np.searchsorted(self.l_grid, l)-1, 0, len(self.l_grid)-2))
        # Interpolazione lineare del guadagno tra due nodi della griglia
        a = (l - self.l_grid[i]) / (self.l_grid[i+1] - self.l_grid[i])
        return (1.0 - a) * self.K_grid[i] + a * self.K_grid[i+1]
\end{lstlisting}

\subsection{Riga 48--58: Fast QP con Algoritmo FISTA}

\begin{theoryframe}
\textbf{Teoria (FISTA):} Per QP con vincoli box
$\min_u \frac{1}{2}u^T H u + g^T u$ s.t. $l_o \le u \le h_i$,
il metodo FISTA (Beck \& Teboulle 2009) applica discesa del gradiente
proiettato con momento di Nesterov:
\[x_{k+1} = \Pi_{[lo,hi]}\!\left(y_k - \frac{Hy_k + g}{L}\right), \quad
y_{k+1} = x_{k+1} + \frac{t_k-1}{t_{k+1}}(x_{k+1}-x_k)\]
con $L = \lambda_{max}(H)$ e $t_{k+1} = \frac{1+\sqrt{1+4t_k^2}}{2}$.
La convergenza e' $O(1/k^2)$ anziche' $O(1/k)$ del gradient descent.
\end{theoryframe}

\begin{lstlisting}[language=Python, caption=solvers.py rr.\ 48-58: FISTA per QP con vincoli box]
def box_qp(H, g, lo, hi, x0=None, iters=60):
    # Stima del passo: 1/L dove L = lambda_max(H)
    L = float(np.linalg.eigvalsh(H)[-1])
    x = np.clip(np.zeros_like(g) if x0 is None else x0, lo, hi)
    y, t = x.copy(), 1.0
    for _ in range(iters):
        # Passo di discesa proiettato sul box [lo, hi]
        x_new = np.clip(y - (H @ y + g) / L, lo, hi)
        # Aggiornamento del momento di Nesterov
        t_new = 0.5 * (1.0 + np.sqrt(1.0 + 4.0 * t * t))
        y = x_new + ((t - 1.0) / t_new) * (x_new - x)
        x, t = x_new, t_new
    return x
\end{lstlisting}

\subsection{Riga 61--76: Condensazione del Problema MPC}

\begin{theoryframe}
\textbf{Teoria (Condensazione):} Date le dinamiche a spazio di stato
$x_{k+1} = Ax_k + Bu_k + c$, sostituendo ricorsivamente si ottiene:
\[\mathbf{X} = \Phi x_0 + \Gamma \mathbf{U} + \mathbf{C}_{st}\]
dove $\mathbf{X} = [x_1; \ldots; x_N]$, $\mathbf{U} = [u_0;\ldots;u_{N-1}]$.
Le matrici $\Phi$, $\Gamma$, $\mathbf{C}_{st}$ si costruiscono iterativamente.
Il costo quadratico si riscrive come QP puro in $\mathbf{U}$:
$H = \Gamma^T S_{blk} \Gamma + W I$, $g = \Gamma^T S_{blk}(\Phi x_0 + C_{st} - X^*)$.
\end{theoryframe}

\begin{lstlisting}[language=Python, caption=solvers.py rr.\ 61-114: Condensazione MPC e classe UpperBodyMPC]
def _condense(A, B, c, N):
    n, m = B.shape
    Phi = np.zeros((N*n, n)); Gam = np.zeros((N*n, N*m)); Cst = np.zeros(N*n)
    Ak = np.eye(n)
    for i in range(N):
        Ak = A @ Ak
        Phi[i*n:(i+1)*n] = Ak
        row = i * n
        prev = Gam[row-n:row] if i > 0 else np.zeros((n, N*m))
        Gam[row:row+n] = A @ prev
        Gam[row:row+n, i*m:(i+1)*m] = B
        Cst[row:row+n] = (A @ Cst[row-n:row] if i > 0 else 0.0) + c
    return Phi, Gam, Cst

class UpperBodyMPC:
    # Disaccoppiamento orizzontale/verticale (Eq. 9 di Cui 2022)
    def _solve(self, A, B, c, x0, x_ref, S, W, lo, hi, warm):
        Phi, Gam, Cst = _condense(A, B, c, self.N)
        Sb = np.kron(np.eye(self.N), S)     # matrice di peso a blocchi
        # Hessiana H e gradiente g del QP condensato
        H = Gam.T @ Sb @ Gam + W * np.eye(self.N)
        g = Gam.T @ Sb @ (Phi @ x0 + Cst - x_ref.reshape(-1))
        warm = np.clip(np.concatenate((warm[1:], warm[-1:])), lo, hi)
        u = box_qp(H, g, lo, hi, warm)
        return u, (Phi @ x0 + Gam @ u + Cst).reshape(self.N, -1)

    def solve(self, x_h, ref_h, h, zdd_ref, ds_max, x_v, ref_v, f_min, f_max):
        dt = self.dt
        # LIP orizzontale: s_ddot = (g+z_ddot)/h * delta_s
        Bh = np.array([[0.5*dt*dt],[dt]]) * (G + zdd_ref) / max(h, 1e-3)
        self.u_h, self.x_h = self._solve(self.Ah, Bh, np.zeros(2),
                               x_h, ref_h, self.S_h, self.W_h, -ds_max, ds_max, self.u_h)
        # Verticale: z_ddot = F_z/m_b - g
        Bv = np.array([[0.5*dt*dt],[dt]]) / self.m_b
        cv = np.array([-0.5*dt*dt*G, -dt*G])   # termine costante della gravita'
        self.u_v, self.x_v = self._solve(self.Ah, Bv, cv,
                               x_v, ref_v, self.S_v, self.W_v, f_min, f_max, self.u_v)
        return float(self.u_h[0]), float(self.u_v[0])
\end{lstlisting}

%=============================================================
\chapter{rotino\_pid: Controllore Baseline PD}
%=============================================================

\section{File: controller.py (PID)}

\subsection{Legge PD di Bilanciamento}

\begin{theoryframe}
\textbf{Teoria (PD sul pendolo inverso):} La legge PD classica per
bilanciare un pendolo inverso su ruote mappa lo stato
$[\theta, \dot\theta, \dot x, x]$ in coppia alle ruote:
$u = K_\theta(\theta-\theta^*) + K_{\dot\theta}\dot\theta
    - K_{\dot x}(\dot x - \dot x^*) - K_x(x - x^*)$.
I guadagni sono determinati per tentativi o tramite linearizzazione
del sistema (VLWIP) e specifica dei poli desiderati.
\end{theoryframe}

\begin{lstlisting}[language=Python, caption=controller.py PID: Legge PD di bilanciamento (righe 644-648)]
k_th, k_thd, k_xd, k_x, max_wheel_state = balance_gains_for_state(self.jump_state)
wheel_u_raw = clamp(
    k_th  * (theta - theta_ref)   # P sull'angolo
    + k_thd * theta_dot            # D sull'angolo
    - k_xd  * vel_err              # D sulla velocita' lineare
    - k_x   * pos_err,             # P sulla posizione
    -max_wheel_state, max_wheel_state)
\end{lstlisting}

\subsection{Servo PD di Posizione per le Gambe}

\begin{lstlisting}[language=Python, caption=controller.py PID: Servo PD hip/knee (righe 454-468)]
def _publish_legs(self, targets, q, qd):
    kps = (HIP_SERVO_KP, HIP_SERVO_KP, KNEE_SERVO_KP, KNEE_SERVO_KP)
    kvs = (HIP_SERVO_KV, HIP_SERVO_KV, KNEE_SERVO_KV, KNEE_SERVO_KV)
    tau = [
        clamp(
            clamp(kp*(target - q[name]), -LEG_TORQUE_MAX, LEG_TORQUE_MAX)  # P
            + clamp(-kv * qd[name], -LEG_DAMPING_MAX, LEG_DAMPING_MAX),    # D (saturato)
            -LEG_TORQUE_MAX, LEG_TORQUE_MAX)
        for kp, kv, target, name in zip(kps, kvs, targets, LEG_JOINTS)]
    self.leg_pub.publish(Float64MultiArray(data=[float(x) for x in tau]))
    return tau
\end{lstlisting}

Il termine D e' saturato separatamente ($\pm$\texttt{LEG\_DAMPING\_MAX}) per
evitare cicli limite bang-bang dovuti a picchi di velocita' angolare durante
transizioni rapide di postura.

%=============================================================
\chapter{rotino\_benchmark: Infrastruttura di Validazione}
%=============================================================

%-------------------------------------------------------------
\section{File: campaign.py}
%-------------------------------------------------------------

\begin{lstlisting}[language=Python, caption=campaign.py: Ciclo di vita di un run (righe 32-68)]
def run_one(law, scenario, duration, out_dir, log_dir):
    kill_stale()   # rimuove residui da run precedenti
    launch = ['ros2', 'launch', f'rotino_{law}', f'rotino_{law}.launch.py',
              'gui:=false', *scenario]
    logger = ['ros2', 'run', 'rotino_benchmark', 'logger', '--ros-args',
              '-p', 'use_sim_time:=true', '-p', f'controller:={law}',
              '-p', f'output_dir:={out_dir}']
    sim = subprocess.Popen(launch, ...)
    time.sleep(6.0)        # attesa avvio Gazebo + controllori
    log = subprocess.Popen(logger, ...)
    time.sleep(duration)   # durata del run
    # Terminazione ordinata: SIGINT poi SIGKILL
    for proc in (log, sim):
        os.killpg(os.getpgid(proc.pid), signal.SIGINT)
    time.sleep(3.0)
    for proc in (log, sim):
        os.killpg(os.getpgid(proc.pid), signal.SIGKILL)
\end{lstlisting}

%-------------------------------------------------------------
\section{File: logger.py}
%-------------------------------------------------------------

\begin{lstlisting}[language=Python, caption=logger.py: Conversione quaternione in angoli Roll/Yaw]
def _imu_cb(self, msg):
    self.imu = msg
    q = msg.orientation
    # Estrazione Roll (rotazione attorno all'asse x) dal quaternione
    sinr_cosp = 2.0 * (q.w * q.x + q.y * q.z)
    cosr_cosp = 1.0 - 2.0 * (q.x*q.x + q.y*q.y)
    self.roll = math.atan2(sinr_cosp, cosr_cosp)
    # Estrazione Yaw (rotazione attorno all'asse z) dal quaternione
    siny_cosp = 2.0 * (q.w * q.z + q.x * q.y)
    cosy_cosp = 1.0 - 2.0 * (q.y*q.y + q.z*q.z)
    self.yaw = math.atan2(siny_cosp, cosy_cosp)
\end{lstlisting}

%-------------------------------------------------------------
\section{File: plot.py -- I 9 Pannelli di Analisi}
%-------------------------------------------------------------

\begin{center}
\begin{tabular}{clp{7cm}}
\toprule
\textbf{N.} & \textbf{File PNG} & \textbf{Contenuto e Rilevanza Teorica} \\
\midrule
1 & tracking\_and\_errors.png  & Errori di tracking di $\theta$, $x$, $z$ per confronto risposta transitoria \\
2 & disturbance\_and\_phase\_portrait.png & Profilo di disturbo e ritratto di fase $(\theta, \dot\theta)$: dimostra la convergenza a 0 dell'orbita \\
3 & smc\_sliding\_surface.png  & Superficie $\sigma(t)$: verifica che il controllore SMC raggiunga e mantenga $\sigma=0$ \\
4 & grf\_and\_joints.png       & Ground Reaction Forces e sforzo dei giunti VMC \\
5 & torque\_distribution.png   & Ripartizione coppia anca/ginocchio e rapporto $|\tau_{hip}|/|\tau_{knee}|$ \\
6 & attitude\_3d.png           & Roll e Yaw vs. tempo: robustezza ai disturbi fuori-piano \\
7 & power\_and\_energy.png     & Potenza meccanica $P = \sum|\tau\omega|$ e energia totale (Joule) \\
8 & horizontal\_phase\_portrait.png & Piano delle fasi $(e_x, \dot e_x)$: dimostra la convergenza all'origine \\
9 & imu\_accelerations.png     & Accelerazioni 3D IMU raw: analisi shock e vibrazioni \\
\bottomrule
\end{tabular}
\end{center}

\begin{lstlisting}[language=Python, caption=plot.py: Calcolo della potenza meccanica totale istantanea]
def plot_power_and_energy(data, out_dir):
    for law in ('mpc', 'smc'):
        d = data[law]
        power = []
        energy = [0.0]
        dt = 0.002   # 500 Hz
        for i in range(len(d['t'])):
            # Potenza ruote: P_wheel = |tau_comune * omega_ruota| * 2 (due ruote)
            pw = abs(d['u'][i] * d['wheel_v'][i]) * 2.0
            # Potenza gambe: P_leg = (|tau_hip * omega_hip| + |tau_knee * omega_knee|) * 2
            pl = (abs(d['hip_u'][i] * d['hip_v'][i])
                  + abs(d['knee_u'][i] * d['knee_v'][i])) * 2.0
            power.append(pw + pl)
            # Integrazione numerica (rettangoli): E(t) = integral_0^t P dt
            energy.append(energy[-1] + (pw + pl) * dt)
\end{lstlisting}

%=============================================================
\chapter{Conclusioni}
%=============================================================

Questo manuale ha dimostrato che ogni riga del codice del progetto RoTino
e' la traduzione diretta di un'equazione matematica della teoria del controllo:

\begin{itemize}
    \item Le \textbf{funzioni di rotazione} in \texttt{kinematics.py} implementano
          la geometria del gruppo $SO(3)$ tramite la formula di Rodrigues e il
          prodotto di rotazioni RPY.
    \item Il \textbf{modello VL-WIP} in \texttt{model.py} estrae automaticamente
          dall'URDF i parametri del modello lineare di Cui et al. (2022),
          garantendo coerenza tra simulazione e modello analitico.
    \item Il \textbf{filtro di Kalman} in \texttt{estimator.py} implementa le
          equazioni (20)-(21) del paper, fondendo IMU e odometria in modo ottimale.
    \item Il \textbf{Super-Twisting} in \texttt{rotino\_smc} garantisce convergenza
          in tempo finito con segnale continuo (no chattering), tramite l'algoritmo
          di Levant con integratore aggiornato a ogni ciclo.
    \item La \textbf{condensazione QP} in \texttt{solvers.py} trasforma un problema
          di programmazione dinamica a orizzonte $N$ in un QP denso di taglia $N$,
          risolvibile in microsecondi con FISTA grazie alla struttura box-constrained.
    \item Il \textbf{benchmark} (\texttt{campaign.py}, \texttt{logger.py},
          \texttt{plot.py}) fornisce 9 metriche di validazione che coprono
          tracking, robustezza ai disturbi, consumo energetico e dinamica di
          assetto tridimensionale.
\end{itemize}

\end{document}
